% GENERATED FILE. Edit canonical lessons, then run npm run generate:books.
% canonical-source-hash: fnv1a64:d1f6936b88198a4c
% canonical-lessons: SA-C437-pravrt, SA-C437-asarah, SA-C437-sikarah, SA-C437-binduh, SA-C437-galati

\chapter{The Rainy Season}
\label{ch:sa-the-rainy-season}
\hlchaptermodality{437}

\begin{chapteropening}
\textbf{By the end of this chapter:} \emph{I can describe the rains, from drizzle to downpour.}

Complete the last lesson of chapter 437: I can describe the rains, from drizzle to downpour.
\end{chapteropening}

\section[\texorpdfstring{\sk{प्रावृट्} (prāvṛṭ)}{प्रावृट् (prāvṛṭ)}]{\sk{प्रावृट्} (prāvṛṭ) — the rainy season}

\label{lesson:SA-C437-pravrt}



\begin{quote}
Before the new one: say the Sanskrit for gets soaked, then the Sanskrit for slippery.
\end{quote}

\subsection*{You'll want to know: \sk{प्रावृट्}}
\textbf{\sk{प्रावृट्}} — \emph{prāvṛṭ} — \textquotedblleft{}the rainy season\textquotedblright{}.

\begin{quote}
\sk{प्रावृट्} \sk{आगच्छति}, \sk{नद्यः} \sk{पूर्णाः} \sk{भवन्ति।}
\end{quote}

\begin{quote}
\emph{prāvṛṭ āgacchati, nadyaḥ pūrṇāḥ bhavanti.}
\end{quote}

\begin{quote}
The rainy season comes, and the rivers fill up.
\end{quote}

A feminine noun.

\begin{cousinweb}
\textbf{\sk{प्रावृट्}} (\emph{prāvṛṭ}) is built from \emph{pra}, \textquotedblleft{}forth\textquotedblright{}, and the root √vṛṣ, \textquotedblleft{}to rain\textquotedblright{}; MW gives \emph{prāvṛṣ}, whose nominative is \emph{prāvṛṭ}, as the rainy season, the rains (MW).
\end{cousinweb}

\begin{grammarlens}[title={What you've built}]
One more word for getting around.
\end{grammarlens}

\subsection*{Your turn}
\begin{itemize}
\raggedright
  \item \emph{Say it:} \emph{prāvṛṭ}
  \item \emph{Say it:} \emph{prāvṛṭ}, once more
  \item \emph{Say it:} \emph{prāvṛṭ āgacchati, nadyaḥ pūrṇāḥ bhavanti.} — The rainy season comes, and the rivers fill up.
  \item \emph{Recall:} say the Sanskrit for gets soaked, then the Sanskrit for slippery, then say \emph{prāvṛṭ} again
\end{itemize}

\begin{tcolorbox}[breakable,colback=teal!4,colframe=teal!35!black,title={Before you move on}]
What does \sk{प्रावृट्} mean? (\textquotedblleft{}The rainy season\textquotedblright{}.) Say it once more.
\end{tcolorbox}
\section[\texorpdfstring{\sk{आसारः} (āsāraḥ)}{आसारः (āsāraḥ)}]{\sk{आसारः} (āsāraḥ) — a downpour}

\label{lesson:SA-C437-asarah}



\begin{quote}
Before the new one: say the Sanskrit for slippery, then the Sanskrit for the rainy season.
\end{quote}

\subsection*{You'll want to know: \sk{आसारः}}
\textbf{\sk{आसारः}} — \emph{āsāraḥ} — \textquotedblleft{}a downpour\textquotedblright{}.

\begin{quote}
\sk{आसारेण} \sk{मार्गाः} \sk{जलेन} \sk{पूर्णाः} \sk{सन्ति।}
\end{quote}

\begin{quote}
\emph{āsāreṇa mārgāḥ jalena pūrṇāḥ santi.}
\end{quote}

\begin{quote}
Because of the downpour the roads are full of water.
\end{quote}

A masculine noun.

\begin{cousinweb}
\textbf{\sk{आसारः}} (\emph{āsāraḥ}) is built from \emph{ā} and the root √sṛ, \textquotedblleft{}to run, to flow\textquotedblright{}; MW gives \emph{āsāra} as a hard shower, a downpour (MW).
\end{cousinweb}

\begin{grammarlens}[title={What you've built}]
One more word for getting around.
\end{grammarlens}

\subsection*{Your turn}
\begin{itemize}
\raggedright
  \item \emph{Say it:} \emph{āsāraḥ}
  \item \emph{Say it:} \emph{āsāraḥ}, once more
  \item \emph{Say it:} \emph{āsāreṇa mārgāḥ jalena pūrṇāḥ santi.} — Because of the downpour the roads are full of water.
  \item \emph{Recall:} say the Sanskrit for slippery, then the Sanskrit for the rainy season, then say \emph{āsāraḥ} again
\end{itemize}

\begin{tcolorbox}[breakable,colback=teal!4,colframe=teal!35!black,title={Before you move on}]
What does \sk{आसारः} mean? (\textquotedblleft{}A downpour\textquotedblright{}.) Say it once more.
\end{tcolorbox}
\section[\texorpdfstring{\sk{शीकरः} (śīkaraḥ)}{शीकरः (śīkaraḥ)}]{\sk{शीकरः} (śīkaraḥ) — drizzle; spray}

\label{lesson:SA-C437-sikarah}



\begin{quote}
Before the new one: say the Sanskrit for the rainy season, then the Sanskrit for a downpour.
\end{quote}

\subsection*{You'll want to know: \sk{शीकरः}}
\textbf{\sk{शीकरः}} — \emph{śīkaraḥ} — \textquotedblleft{}drizzle; spray\textquotedblright{}.

\begin{quote}
\sk{प्रातः} \sk{शीकराः} \sk{पतन्ति।}
\end{quote}

\begin{quote}
\emph{prātaḥ śīkarāḥ patanti.}
\end{quote}

\begin{quote}
In the morning a drizzle falls.
\end{quote}

A masculine noun.

\begin{cousinweb}
\textbf{\sk{शीकरः}} (\emph{śīkaraḥ}) is the word MW gives, mostly in the plural, for fine or drizzling rain and spray; MW gives no derivation, so its origin is uncertain (MW).
\end{cousinweb}

\begin{grammarlens}[title={What you've built}]
One more word for getting around.
\end{grammarlens}

\subsection*{Your turn}
\begin{itemize}
\raggedright
  \item \emph{Say it:} \emph{śīkaraḥ}
  \item \emph{Say it:} \emph{śīkaraḥ}, once more
  \item \emph{Say it:} \emph{prātaḥ śīkarāḥ patanti.} — In the morning a drizzle falls.
  \item \emph{Recall:} say the Sanskrit for the rainy season, then the Sanskrit for a downpour, then say \emph{śīkaraḥ} again
\end{itemize}

\begin{tcolorbox}[breakable,colback=teal!4,colframe=teal!35!black,title={Before you move on}]
What does \sk{शीकरः} mean? (\textquotedblleft{}Drizzle; spray\textquotedblright{}.) Say it once more.
\end{tcolorbox}
\section[\texorpdfstring{\sk{बिन्दुः} (binduḥ)}{बिन्दुः (binduḥ)}]{\sk{बिन्दुः} (binduḥ) — a drop}

\label{lesson:SA-C437-binduh}



\begin{quote}
Before the new one: say the Sanskrit for a downpour, then the Sanskrit for drizzle.
\end{quote}

\subsection*{You'll want to know: \sk{बिन्दुः}}
\textbf{\sk{बिन्दुः}} — \emph{binduḥ} — \textquotedblleft{}a drop\textquotedblright{}.

\begin{quote}
\sk{पत्रे} \sk{जलस्य} \sk{बिन्दुः} \sk{अस्ति।}
\end{quote}

\begin{quote}
\emph{patre jalasya binduḥ asti.}
\end{quote}

\begin{quote}
There is a drop of water on the leaf.
\end{quote}

A masculine noun.

\begin{cousinweb}
\textbf{\sk{बिन्दुः}} (\emph{binduḥ}) is the word MW gives for a drop, a dot, a spot; MW gives no derivation, so its origin is uncertain; CDIAL 9240 records its later forms (MW; CDIAL 9240).
\end{cousinweb}

\begin{grammarlens}[title={What you've built}]
One more word for getting around.
\end{grammarlens}

\subsection*{Your turn}
\begin{itemize}
\raggedright
  \item \emph{Say it:} \emph{binduḥ}
  \item \emph{Say it:} \emph{binduḥ}, once more
  \item \emph{Say it:} \emph{patre jalasya binduḥ asti.} — There is a drop of water on the leaf.
  \item \emph{Recall:} say the Sanskrit for a downpour, then the Sanskrit for drizzle, then say \emph{binduḥ} again
\end{itemize}

\begin{tcolorbox}[breakable,colback=teal!4,colframe=teal!35!black,title={Before you move on}]
What does \sk{बिन्दुः} mean? (\textquotedblleft{}A drop\textquotedblright{}.) Say it once more.
\end{tcolorbox}
\section[\texorpdfstring{\sk{गलति} (galati)}{गलति (galati)}]{\sk{गलति} (galati) — drips; trickles}

\label{lesson:SA-C437-galati}



\begin{quote}
Before the new one: say the Sanskrit for drizzle, then the Sanskrit for a drop.
\end{quote}

\subsection*{You'll want to know: \sk{गलति}}
\textbf{\sk{गलति}} — \emph{galati} — \textquotedblleft{}drips; trickles\textquotedblright{}.

\begin{quote}
\sk{पत्रात्} \sk{जलं} \sk{गलति।}
\end{quote}

\begin{quote}
\emph{patrāt jalaṃ galati.}
\end{quote}

\begin{quote}
Water drips from the leaf.
\end{quote}

Past (it …): \sk{जलं} \sk{पत्रात्} \sk{गलितवत्}, \emph{jalaṃ patrāt galitavat}.

\begin{cousinweb}
\textbf{\sk{गलति}} (\emph{galati}) comes from the root √gal, \textquotedblleft{}to drip, to ooze\textquotedblright{}; MW gives \emph{galati} as to drip, drop, trickle, with the participle \emph{galita} (MW).
\end{cousinweb}

\begin{grammarlens}[title={What you've built}]
One more word for getting around.
\end{grammarlens}

\subsection*{Your turn}
\begin{itemize}
\raggedright
  \item \emph{Say it:} \emph{galati}
  \item \emph{Say it:} \emph{galati}, once more
  \item \emph{Say it:} \emph{patrāt jalaṃ galati.} — Water drips from the leaf.
  \item \emph{Recall:} say the Sanskrit for drizzle, then the Sanskrit for a drop, then say \emph{galati} again
\end{itemize}

\begin{tcolorbox}[breakable,colback=teal!4,colframe=teal!35!black,title={Before you move on}]
What does \sk{गलति} mean? (\textquotedblleft{}Drips; trickles\textquotedblright{}.) Say it once more.
\end{tcolorbox}
